\documentclass[11pt]{article}
\usepackage[margin=1in]{geometry}
\begin{document}

\begin{flushright}
Yaoping~Wang \\
Department of Statistics, Rutgers University \\
yaoping.wang@rutgers.edu \\
\today
\end{flushright}

\vspace{2mm}

Dear Editor,

I submit the manuscript ``A comparison index for functional ranking in model selection'' for consideration in {\it Biometrika}.  This work is original and not under consideration elsewhere.

The paper introduces a comparison index $\kappa$ that governs when a scalar selection functional (AIC, CP width, cross-validated error) preserves the risk ranking of candidate estimators.  The main result is a phase transition: when the population ordering agrees, misalignment decays at rate $\exp(-n\kappa^2/2)$; otherwise it persists asymptotically.  A finite-sample consequence is the {\it competition window}, a non-monotonic misalignment trajectory at intermediate sample sizes that provides a testable diagnostic distinguishing competition from structural bias-crossing.  Because $\kappa$ can be estimated directly from calibration data without cross-validation, practitioners have a computable signal-to-noise measure for deciding when a given selection criterion can be trusted.

Recent work on conformal prediction after selection (Hegazy et al., 2025; Liang et al., 2026; Yang \& Kuchibhotla, 2024) focuses on maintaining coverage validity post-selection.  The present paper addresses a distinct question---whether the selection functional can distinguish competing estimators in the first place---and characterises its fundamental resolution limit.  The work bridges Le Cam's comparison of statistical experiments with modern model selection practice, connecting classical asymptotic decision theory to emerging problems in conformal inference, a direction that is well within the methodological remit of this journal.

\vspace{3mm}
{\bf Declaration of the use of generative AI and AI-assisted technologies:}
During the preparation of this work the author used Claude Code (Anthropic) in order to assist with coding, editing and literature review.  After using this tool the author reviewed and edited the content as necessary and takes full responsibility for the content of the publication.

\vspace{5mm}

Yours sincerely,

Yaoping Wang

\end{document}
